\section{Related Work}
\label{sec:related-work}

\emph{Approximate-FHE attacks and defenses.} Li and Micciancio identified the secret-dependent approxi\-mation-error channel in CKKS-style decryption~\cite{lm}.
LMSS introduced the IND-CPAD formulation and used differential-privacy and noise-flooding techniques to recover security under explicit parameter conditions~\cite{lmss}.
Subsequent work has sharpened the practical precision tradeoff~\cite{costache2023precision} and revisited concrete noise-flooding parameters~\cite{bergamaschi2025revisiting}.
Attacks against exact FHE at aggressive parameters show that decryption-query security is not confined to approximate schemes~\cite{cheon2024attacks}.
Our contribution is neither a new attack nor a better concrete parameter set: it machine-checks an existing reduction and the square-root composition mechanism on which such parameters rely.

\emph{Mechanized game-based cryptography.} SSProve combines state-separating package algebra with a probabilistic relational program logic in Rocq and validates the framework on encryption, KEM--DEM, and sigma-protocol case studies~\cite{ssprove}.
Its central bridge lifts invariant-preserving exact judgments for individual procedures to perfect indistinguishability of packages.
Our result follows the same security-facing pattern at a different quantitative boundary: it lifts a conditional-KL judgment for one oracle implementation to additive indistinguishability of a complete adaptive program.
We retain SSProve's syntax, package algebra, heaps, linking, and subdistribution semantics, while adding judgments and a selected-call compiler.
Recent Nominal-SSProve work makes package state separation intrinsic and reports errors found in an ElGamal mechanization~\cite{nominal-ssprove}; a further development mechanizes query-counting, multi-instance, and nested hybrid arguments~\cite{nested-hybrids}.
Those results address exact package composition and general hybrid structure.
Our compiler instead preserves an adaptive sequence of conditional KL costs and derives a sublinear statistical-distance loss.

EasyCrypt is another popular framework for game-based cryptographic proofs~\cite{easycrypt}.
Adversaries are treated as opaque modules, and reflection techniques~\cite{reflection} can be used to extract the adversary's semantics.
Our abstraction boundary differs because SSProve adversaries are concrete program syntax: the verified trace compiler factors selected calls before the replacement theorem is applied.
This compiler supplies a checked connection between local divergence bounds and complete adaptive SSProve games.
EasyCrypt also has a cost logic for bounding execution costs and oracle-call counts of adversarial programs~\cite{easycrypt-cost}.

Squirrel is a standalone interactive prover based on a higher-order extension of the Computationally Complete Symbolic Attacker (CCSA) logic.
Baelde et al. extend this logic with explicit failure-probability and computational-indistinguishability bounds; the latter depend on adversary execution time and oracle-call counts under a Turing-machine cost model~\cite{squirrel-concrete}.
Their proof transformations avoid exponential losses caused by duplicated induction hypotheses.
Fontaine et al. implement bound-annotated tactics and proof-term elaboration for the reachability fragment, leaving concrete indistinguishability with precise complexity bounds outside the scope of their implementation~\cite{squirrel-errors}.
This reachability reasoning shares additive probability-error accounting with our \(\AEJ\) judgment, but its bounds concern formula-failure probabilities over random tapes; \(\AEJ\) instead bounds postcondition-failure probability in an existential coupling of completed program outputs.

\emph{Program logics for differential privacy.}
Approximate relational Hoare logic (apRHL) of \cite{certipriv} and its follow up work fpRHL \cite{aprhl} are relational program logics for cryptographic games in differential privacy settings.

apRHL was initially formalized in Coq as part of the CertiPriv framework,
and was later integrated into EasyCrypt as well \cite{aprhl-couplings, aprhl-advanced}.
The settings it targets are incomparable with those of homomorphic encryptions.
In particular, their security bounds allow multiplicative errors,
but still scales linearly in the number of oracle calls.

fpRHL \emph{does} allow sub-linear compositions much like our Pythagorean logic.
The logic is parameterized by $f$-divergences, which is a more general notion than KL-divergence.
To the best of our knowledge, fpRHL has not been formalized or verified in any proof assistant software.
For our purpose, it would have been a valid strategy to implement fpRHL instead,
or its specialization to KL-divergences.
From there, it still remains to treat adaptive adversaries,
which would require either extending fpRHL with EasyCrypt-style ``adversarial call'' rules,
or using our compiler again.
On the other hand, we wanted to model Micciancio--Walter Pythagorean preservation~\cite{mw17} itself as an inference rule,
and to verify the soundness of any program logic rules in SSProve.
The rest of our design decisions followed accordingly.
